\chapter{Jets, Infrared-Safe Observables, and Hadronization}
\label{ch:jets-ir-safe-observables-hadronization}
\ConventionsMixed

Jets enter quantum field theory as observables built from the energy and
momentum deposited by physical final states.  The final states in QCD are
hadronic states in the physical Hilbert space; quarks and gluons are
short-distance variables used inside a gauge-fixed or factorized calculation.
The mathematical question is therefore not how to assign reality to colored
partons, but when a measurement on physical final states admits a controlled
short-distance expansion in which partonic variables are legitimate internal
coordinates.

\section{Weighted Cross Sections and Measurement Functions}

Let \(\ket{i}\) be an incoming scattering state and let \(X\) range over a
complete family of outgoing physical states.  A detector observable is a
measurable function \(F(X)\) of the final state.  Its weighted inclusive cross
section has the form
\begin{equation}
  \sigma[F]
  =
  \frac{1}{\mathcal N_i}
  \sum_X
  (2\pi)^D\delta^{(D)}(P_i-P_X)
  |\langle X|T|i\rangle|^2\,F(X),
  \label{eq:jet-weighted-cross-section-hilbert}
\end{equation}
where \(\mathcal N_i\) is the incoming flux normalization and \(T\) is the
transition operator defined by \(S=\mathbf 1+\ii T\).  Equation
\eqref{eq:jet-weighted-cross-section-hilbert} is stated in terms of physical
states.  A perturbative partonic formula is a further expansion of the same
quantity, after an infrared regulator and an ultraviolet renormalization
scheme have been chosen:
\begin{equation}
  \sigma_{\rm part}[F]
  =
  \sum_{n\geq0}
  \int \dd\Phi_n\,
  |\mathcal M_n(p_1,\ldots,p_n)|^2\,
  F_n(p_1,\ldots,p_n).
  \label{eq:jet-partonic-measurement-function}
\end{equation}
Here \(\dd\Phi_n\) is the \(n\)-body Lorentz-invariant phase-space measure,
\(\mathcal M_n\) is a renormalized partonic matrix element, and \(F_n\) is the
partonic representative of the detector measurement.  The functions \(F_n\)
must satisfy compatibility conditions at the boundaries of phase space where a
parton becomes soft or two partons become collinear.

\begin{definition}[Infrared and collinear safety]
\label{def:irc-safe-measurement}
A family of measurement functions \(F=\{F_n\}_{n\geq0}\) is infrared and
collinear safe, abbreviated IRC safe, if the following limits hold on the
phase-space regions relevant to the observable:
\begin{align}
  \lim_{\lambda\downarrow0}
  F_{n+1}(p_1,\ldots,p_n,\lambda q)
  &=
  F_n(p_1,\ldots,p_n),
  \label{eq:irc-soft-safety}\\
  \lim_{\theta\downarrow0}
  F_{n+1}(p_1,\ldots,zp,(1-z)p,\ldots,p_n)
  &=
  F_n(p_1,\ldots,p,\ldots,p_n),
  \qquad 0<z<1.
  \label{eq:irc-collinear-safety}
\end{align}
The convergence is required with enough uniformity that the difference between
the two sides is integrable against the soft and collinear factorization
kernels of the theory at the perturbative order under consideration.
Equivalently, near a single unresolved boundary and away from other unresolved
boundaries, there must be constants \(C,\alpha>0\), depending on the compact
set of resolved momenta but not on the unresolved scale, such that
\begin{align}
  |F_{n+1}(p_1,\ldots,p_n,\lambda q)-F_n(p_1,\ldots,p_n)|
  &\le C\lambda^\alpha,\\
  |F_{n+1}(p_1,\ldots,zp,(1-z)p,\ldots,p_n)
    -F_n(p_1,\ldots,p,\ldots,p_n)|
  &\le C\theta^\alpha ,
\end{align}
where \(\theta\) is the splitting angle.  On corners where several emissions
are unresolved, the same estimate is required after replacing each unresolved
cluster by its reduced soft or collinear configuration.  These estimates make
the products with \(\dd\lambda/\lambda\) and
\(\dd\theta^2/\theta^2\) locally integrable after the real--virtual
subtraction.
\end{definition}

\begin{proposition}[Fixed-order finiteness criterion for IRC-safe measurements]
\label{prop:irc-safety-fixed-order-finiteness}
Consider an inclusive observable in a massless gauge theory for which the
renormalized real and virtual matrix elements admit the soft and collinear
boundary factorizations displayed below, with remainders locally integrable
after subtracting the reduced lower-multiplicity channel.  If the measurement
family \(F\) is IRC safe in the sense of
Definition~\ref{def:irc-safe-measurement}, then real-emission singularities
and virtual singularities cancel locally on the unresolved phase-space
boundaries after the KLN sum over degenerate final states.  With incoming
hadrons, the remaining initial-state collinear singularities are absorbed into
the renormalized parton distribution functions.
\end{proposition}

\begin{proof}
In a soft region, with an extra gluon momentum \(k=\lambda q\), the real
matrix element factorizes as
\[
  |\mathcal M_{n+1}|^2\,\dd\Phi_{n+1}
  =
  |\mathcal M_n|^2\,\dd\Phi_n\,
  \left[
    \sum_{i,j}\mathcal S_{ij}(q)
  \right]
  \frac{\dd\lambda}{\lambda}\dd\Omega_q
  +R_{\rm soft}.
\]
The remainder \(R_{\rm soft}\) is locally integrable by the hypothesis of the
proposition.
The eikonal factor \(\mathcal S_{ij}\) is the same factor that appears with
opposite sign in the virtual soft correction to the \(n\)-parton channel.
Equation \eqref{eq:irc-soft-safety} replaces the real measurement by
\(F_n\) on the singular part of the integral, so the real and virtual soft
integrals have the same weight and cancel.

In a final-state collinear region, two momenta approach
\[
  p_a=zp+r,\qquad p_b=(1-z)p-r,
  \qquad r_\perp^2\to0.
\]
The real matrix element and phase space factorize as
\[
  |\mathcal M_{n+1}|^2\,\dd\Phi_{n+1}
  =
  |\mathcal M_n(\ldots,p,\ldots)|^2\,\dd\Phi_n\,
  \frac{\alpha_s}{2\pi}
  P_{a b\leftarrow c}(z)
  \frac{\dd r_\perp^2}{r_\perp^2}\dd z
  +R_{\rm coll}.
\]
The remainder \(R_{\rm coll}\) is locally integrable by the hypothesis of the
proposition.
Equation \eqref{eq:irc-collinear-safety} gives the same measurement weight to
the unresolved pair and to its parent momentum.  The singular integral then
cancels against the corresponding virtual collinear correction in the KLN
sum.  If the collinear parton is in the initial state of a hadron collision,
the same factorization kernel multiplies the operator renormalization of the
parton distribution function rather than canceling inside the hard partonic
subprocess.  The finite hard coefficient is obtained after this mass
factorization step.
\end{proof}

Examples of IRC-safe observables include the total \(e^+e^-\) hadronic cross
section, thrust
\[
  T(X)
  =
  \max_{\widehat n}
  \frac{\sum_{h\in X}|\vec p_h\cdot\widehat n|}
       {\sum_{h\in X}|\vec p_h|},
\]
the \(C\)-parameter, jet masses defined with an IRC-safe algorithm, and the
energy correlators of
Chapter~\ref{ch:volume-ii-19}.  Observables that measure an identified
hadron, a charged-particle multiplicity, or a single unsummed collinear
fragment require additional nonperturbative input such as fragmentation
functions or detector-specific response functions.

\section{Sterman--Weinberg Jets}

The original two-jet observable of Sterman and Weinberg is a clean example of
the preceding definition.  In \(e^+e^-\) annihilation at center-of-mass energy
\(Q\), choose two parameters
\[
  0<\epsilon_{\rm out}<1,\qquad 0<\delta_{\rm cone}<\pi.
\]
An event is counted as a two-jet event if there exist two cones of half-angle
\(\delta_{\rm cone}\) such that the total energy outside the two cones is at
most \(\epsilon_{\rm out}Q\).  The corresponding measurement function is
\[
  F_{\rm SW}(X)
  =
  \begin{cases}
  1,& \text{if }X\text{ satisfies the two-cone energy condition},\\
  0,& \text{otherwise}.
  \end{cases}
\]

Soft safety follows because adding a particle of energy \(\lambda Q\) changes
the out-of-cone energy by at most \(\lambda Q\).  Away from the boundary
where the out-of-cone energy is exactly \(\epsilon_{\rm out}Q\), the
classification is unchanged for sufficiently small \(\lambda\).  The boundary
has zero measure in the phase-space integral, and the unresolved soft
singularity is integrated with the same measurement weight as the event
without the soft particle.

Collinear safety follows because a sufficiently collinear splitting of a
momentum inside one of the cones remains inside the same cone, while a
splitting outside the cones preserves the total out-of-cone energy in the
collinear limit.  The singular collinear phase-space integral therefore sees
the parent momentum measurement.

For very small \(\epsilon_{\rm out}\) and \(\delta_{\rm cone}\), the fixed-order
series contains logarithms such as
\[
  \alpha_s^m
  \log^r\epsilon_{\rm out}\,
  \log^s\delta_{\rm cone},
  \qquad r+s\leq 2m.
\]
The two-jet cross section is then still an IRC-safe observable, but the useful
short-distance expansion is a resummed expansion rather than a fixed-order
polynomial in \(\alpha_s(Q)\).

\section{Sequential Recombination Jet Algorithms}

Modern collider jet definitions are usually algorithmic.  For a hadron
collider, let \(p_{Ti}\) be transverse momentum, \(y_i\) rapidity, and
\(\phi_i\) azimuth.  Define
\[
  \Delta R_{ij}^2=(y_i-y_j)^2+(\phi_i-\phi_j)^2.
\]
The generalized \(k_T\) family with radius \(R\) and exponent \(p\) uses
\begin{equation}
  d_{ij}
  =
  \min(p_{Ti}^{2p},p_{Tj}^{2p})
  \frac{\Delta R_{ij}^2}{R^2},
  \qquad
  d_{iB}=p_{Ti}^{2p}.
  \label{eq:generalized-kt-distance}
\end{equation}
At each step the smallest number among all \(d_{ij}\) and \(d_{iB}\) is
selected.  If it is \(d_{ij}\), the two objects are recombined, usually by
adding their four-momenta.  If it is \(d_{iB}\), object \(i\) is declared a
jet and removed from the list.  The cases
\[
  p=1,\qquad p=0,\qquad p=-1
\]
give the \(k_T\), Cambridge--Aachen, and anti-\(k_T\) algorithms,
respectively.

\begin{figure}[H]
  \centering
  \resizebox{0.92\textwidth}{!}{%
  \begin{tikzpicture}[x=1cm,y=1cm,>=Stealth,every node/.style={font=\scriptsize}]
    \node[draw,rounded corners=2pt,align=center,minimum width=3.2cm,minimum height=0.9cm]
      (state) at (0,0) {physical final state\\hadrons or calorimeter cells};
    \node[draw,rounded corners=2pt,align=center,minimum width=3.0cm,minimum height=0.9cm]
      (measure) at (4.0,0) {measurement\\\(F(X)\)};
    \node[draw,rounded corners=2pt,align=center,minimum width=3.3cm,minimum height=0.9cm]
      (irc) at (8.0,0) {soft and collinear\\compatibility};
    \node[draw,rounded corners=2pt,align=center,minimum width=3.3cm,minimum height=0.9cm]
      (calc) at (12.1,0) {fixed order,\\factorization, resummation};
    \draw[->,thick] (state)--(measure);
    \draw[->,thick] (measure)--(irc);
    \draw[->,thick] (irc)--(calc);

    \node[align=center,blue!65!black] at (4.0,-1.2)
      {jet algorithm, event shape,\\energy correlator};
    \node[align=center,purple!70!black] at (8.0,-1.2)
      {\(F_{n+1}\to F_n\) on\\unresolved boundaries};
    \node[align=center,orange!80!black] at (12.1,-1.2)
      {partonic variables are\\internal coordinates};
  \end{tikzpicture}
  }
  \caption{A perturbative jet calculation begins with a physical measurement.
  IRC safety is the compatibility condition that lets the measurement be
  represented on unresolved partonic phase-space boundaries.}
  \label{fig:jets-measurement-irc-flow}
\end{figure}

\begin{proposition}[IRC safety of generalized \(k_T\) clustering]
\label{prop:generalized-kt-irc-safety}
For \(p=1,0,-1\), the generalized \(k_T\) clustering algorithm with
four-momentum recombination defines IRC-safe jet momenta and IRC-safe
cross sections after imposing a finite jet transverse-momentum threshold.
\end{proposition}

\begin{proof}
For a collinear splitting \(p\to zp,(1-z)p\), the angular distance satisfies
\(\Delta R_{ab}\to0\), so \(d_{ab}\to0\).  The two collinear objects are
therefore recombined before they can affect finite-angle clustering, and
four-momentum recombination restores the parent momentum.

For a soft momentum \(q\), the effect of recombination with a hard object
changes the hard jet momentum by \(O(q)\).  If the soft object is clustered
among other soft objects or declared as a separate jet, a finite jet
transverse-momentum threshold removes it in the soft limit.  In the anti-\(k_T\)
case, a soft object within radius \(R\) of a hard jet has pair distance
controlled by the hard transverse momentum, so it is pulled into the hard jet
without moving the hard boundary at leading power.  The \(k_T\) and
Cambridge--Aachen cases may cluster soft radiation differently, but the
resulting change of any finite-energy jet momentum vanishes with the soft
energy.  Thus the measurement functions associated with the algorithm obey
\eqref{eq:irc-soft-safety} and \eqref{eq:irc-collinear-safety}.
\end{proof}

The anti-\(k_T\) algorithm is widely used because hard jets acquire stable,
approximately circular catchment areas in the rapidity--azimuth plane while
retaining IRC safety.  Cone algorithms require more care.  A seedless
infrared-safe cone algorithm can be defined by searching over all stable
cones.  A seed-based cone algorithm with a transverse-momentum seed threshold
can fail soft safety: adding an arbitrarily soft seed may create a stable cone
that was absent before, changing the hard jets by a finite amount.

\section{Factorization and Resummation}

IRC safety makes fixed-order perturbation theory finite, but it does not make
every region of the observable distribution well approximated by a few fixed
orders.  When an event shape \(e(X)\) is small, the final state is constrained
to contain narrow energetic jets plus soft radiation.  A factorized
description separates the scales.  For a two-jet event shape in
\(e^+e^-\), a typical factorized form is
\begin{equation}
  \frac{\dd\sigma}{\dd e}
  =
  H(Q,\mu)
  \int \dd s_1\,\dd s_2\,\dd k\,
  J(s_1,\mu)J(s_2,\mu)S(k,\mu)\,
  \delta\!\left(e-\frac{s_1+s_2}{Q^2}-\frac{k}{Q}\right)
  +
  O(e^0).
  \label{eq:jet-event-shape-factorization}
\end{equation}
The hard function \(H\) contains fluctuations at virtuality \(Q^2\), the jet
functions \(J\) contain collinear fluctuations at smaller invariant mass, and
the soft function \(S\) is a vacuum matrix element of Wilson lines measuring
soft energy flow.  Their renormalization-group equations resum logarithms of
the separated scales.

Soft-collinear effective theory gives a systematic operator language for
factorizations of the form
\eqref{eq:jet-event-shape-factorization}.  Its assumptions are visible in the
operator definitions: mode scaling, Wilson-line decoupling at leading power,
measurement factorization, and a power expansion in the small event-shape
parameter.  The same logic applies to hadron-collider observables, with the
additional parton distribution functions needed for incoming hadrons and with
factorization-breaking channels treated as separate hypotheses.

\subsection{Mode Expansion and Logarithmic Accuracy}

For a two-jet observable with small dimensionless event shape \(e\), choose
lightlike vectors \(n,\bar n\) with \(n\cdot\bar n=2\) and write momenta as
\[
  p^\mu=\frac{n\cdot p}{2}\bar n^\mu
       +\frac{\bar n\cdot p}{2}n^\mu+p_\perp^\mu .
\]
The SCET expansion parameter \(\lambda\) is fixed by the measured region.  For
thrust-like two-jet kinematics, \(e\sim\lambda^2\), and the momentum scalings
are
\[
  p_n^\mu\sim Q(\lambda^2,1,\lambda),\qquad
  p_{\bar n}^\mu\sim Q(1,\lambda^2,\lambda),\qquad
  p_s^\mu\sim Q(\lambda^2,\lambda^2,\lambda^2),
\]
where the three entries denote
\((n\cdot p,\bar n\cdot p,p_\perp)\).  A leading-power factorization theorem
is a statement that, for \(e\ll1\) and away from additional measured
hierarchies,
\[
  \frac{\dd\sigma}{\dd e}
  =
  H\otimes J_n\otimes J_{\bar n}\otimes S
  +
  O(\lambda)
  \quad\text{or}\quad
  O(e^{1/2}),
\]
with the precise power depending on the measurement and recoil convention.
The symbol \(\otimes\) denotes the convolution variables in the displayed
factorization formula, not an abstract tensor product.  The hard function is a
Wilson coefficient of the QCD current matched onto SCET operators; the jet
functions are matrix elements of collinear fields with collinear Wilson
lines; the soft function is a vacuum matrix element of soft Wilson lines with
the measurement operator inserted.

\begin{hypothesis}[Leading-power SCET factorization datum]
\label{hyp:leading-power-scet-factorization}
A leading-power SCET factorization statement for an observable \(F\) consists
of a small parameter \(\lambda\), a list of modes with momentum scaling, a
mode-separated operator basis through the declared order in \(\lambda\), a
measurement decomposition on those modes, a statement that Glauber or other
factorization-breaking exchanges either cancel or are represented by
additional operators, and a remainder estimate in powers of \(\lambda\) on the
kinematic region where the formula is used.
\end{hypothesis}

Renormalization of the hard, jet, and soft functions gives evolution
equations.  If \(L=\log e\), then logarithmic accuracy means a definite
truncation of the exponent and matching functions.  A prediction at LL
retains all terms \(\alpha_s^mL^{m+1}\) in the cumulant exponent generated by
one-loop cusp evolution and tree matching.  NLL adds the next tower
\(\alpha_s^mL^m\), requiring higher cusp data, one-loop non-cusp anomalous
dimensions, and running-coupling effects.  NNLL adds
\(\alpha_s^mL^{m-1}\), with the corresponding higher anomalous dimensions and
matching coefficients.  The labels LL, NLL, and NNLL therefore refer to a
specified observable, factorization theorem, scale choice, and matching
scheme; they do not by themselves state a numerical error bound outside the
declared logarithmic hierarchy.

\begin{controlledapproximation}[Perturbative prediction for an IRC-safe jet observable]
A perturbative prediction for a jet observable consists of:
the renormalized fixed-order hard functions through a stated order in
\(\alpha_s\); the logarithmic accuracy of any resummation, such as LL, NLL,
NNLL, or beyond; the factorization theorem and its power corrections; the
renormalized PDFs or fragmentation functions when they enter; and a
prescription for matching resummed and fixed-order regions without double
counting.
\end{controlledapproximation}

\section{Jet Substructure Observables and Grooming}

Jet substructure observables are measurements on the constituents of a
previously identified jet.  The jet definition, recombination scheme, grooming
algorithm, and substructure measurement are all part of the observable.  Their
IRC behavior must therefore be checked as a combined map from the final-state
calorimetric measure to a number or distribution.

\begin{definition}[Soft drop]
\label{def:soft-drop-grooming}
Start with the constituents of a jet of radius \(R_0\), recluster them with
the Cambridge--Aachen metric, and recursively decluster the last clustering
step into two branches \(i,j\).  Let
\[
  z_{ij}=\frac{\min(p_{Ti},p_{Tj})}{p_{Ti}+p_{Tj}},
  \qquad
  \theta_{ij}=\frac{\Delta R_{ij}}{R_0}.
\]
For parameters \(z_{\rm cut}>0\) and \(\beta_{\rm SD}\geq0\), the soft-drop
condition is
\[
  z_{ij}>z_{\rm cut}\,\theta_{ij}^{\beta_{\rm SD}} .
\]
If the condition fails, the softer branch is removed and the recursion
continues on the harder branch.  If the condition holds, the remaining
constituents define the groomed jet.  The case \(\beta_{\rm SD}=0\) is the
modified mass-drop tagger.
\end{definition}

\begin{proposition}[IRC behavior of soft drop for \(\beta_{\rm SD}\geq0\)]
\label{prop:soft-drop-irc-safety}
With an IRC-safe input jet algorithm and a finite transverse-momentum
threshold for the final groomed jet, soft drop with
\(\beta_{\rm SD}\geq0\) gives IRC-safe groomed four-momentum and IRC-safe
groomed jet mass away from the boundary hypersurfaces of the measurement.
\end{proposition}

\begin{proof}
Add a particle of transverse momentum \(\lambda p_T\) with
\(\lambda\downarrow0\).  If it is tested against a hard branch at finite
angle, then \(z_{ij}=O(\lambda)\), while
\(z_{\rm cut}\theta_{ij}^{\beta_{\rm SD}}\) has a positive lower bound on
compact subsets away from \(\theta_{ij}=0\).  The soft branch is removed and
cannot change the groomed four-momentum except on the codimension-one
surface where the inequality is exactly saturated.  If the soft particle is
collinear to a hard branch, Cambridge--Aachen recombines or declusters it at
small \(\theta\), and its contribution to the branch momentum vanishes with
\(\lambda\).

For a collinear splitting of a hard particle, the two branches have
\(\theta_{ij}\downarrow0\).  When \(\beta_{\rm SD}\geq0\), the right side
\(z_{\rm cut}\theta_{ij}^{\beta_{\rm SD}}\) tends to \(0\) for
\(\beta_{\rm SD}>0\) and remains \(z_{\rm cut}\) for
\(\beta_{\rm SD}=0\).  In the first case, any fixed energy fraction passes the
condition for sufficiently small angle and the four-momentum recombines to
the parent.  In the second case, a sufficiently asymmetric collinear split may
remove the soft-collinear branch, but the removed momentum is collinear and
has vanishing contribution to the groomed mass in the collinear limit; the
groomed four-momentum changes by the removed energy fraction only on the
measurement boundary relevant to flavor-sensitive or track-sensitive
variants.  For the stated groomed mass and four-momentum observables, the
measurement has the IRC limits of
Definition~\ref{def:irc-safe-measurement}.
\end{proof}

\begin{definition}[\(N\)-subjettiness]
\label{def:n-subjettiness}
For a jet with constituents \(k\), choose \(N\) axes
\(\{a_1,\ldots,a_N\}\) by a specified IRC-safe axis prescription.  For
\(\beta_\tau>0\),
\[
  \tau_N^{(\beta_\tau)}
  =
  \frac{1}{p_{TJ}R_0^{\beta_\tau}}
  \sum_{k\in J}p_{Tk}\,
  \min_{a=1,\ldots,N}(\Delta R_{ka})^{\beta_\tau}.
\]
Ratios such as \(\tau_{21}=\tau_2/\tau_1\) are substructure coordinates only
on the domain where the denominator is separated from zero; near the
denominator boundary their distribution requires the same factorization and
resummation data as the underlying \(\tau_N\).
\end{definition}

Energy-correlation functions, defined in
Chapter~\ref{ch:volume-ii-19}, provide a detector-based substructure
alternative.  They depend only on the positive calorimetric measure and
angular kernels, whereas \(N\)-subjettiness also depends on an axis-finding
map.  Both are useful, but their mathematical inputs differ.

\section{Shape Functions and Track Functions}

When a measured event shape is of order \(\Lambda_{\rm QCD}/Q\), the soft
function in \eqref{eq:jet-event-shape-factorization} contains
nonperturbative matrix elements.  A typical leading shape-function
factorization has the form
\[
  \frac{\dd\sigma}{\dd e}
  =
  \int_0^\infty\dd\ell\,
  \frac{\dd\sigma_{\rm pert}}{\dd e}
  \left(e-\frac{\ell}{Q}\right)
  F_{\rm shape}(\ell)
  +
  O(Q^{-2})
\]
on the region where a single soft momentum variable \(\ell\) controls the
leading power correction.  The function \(F_{\rm shape}\) is a
renormalized matrix element of Wilson lines with a soft measurement operator,
normalized in the chosen scheme.  A fitted parametric form for
\(F_{\rm shape}\) is phenomenological input; the operator definition and its
renormalization are the QFT data.

Track-based observables measure only charged hadrons.  They are generally not
IRC safe as ordinary partonic measurements, because a collinear splitting can
change the charged energy fraction observed by the detector.  The replacement
input is a track function.  For a parton species \(i\), let \(T_i(x,\mu)\) be
the renormalized probability density, in a specified operator scheme, that
the charged tracks in the jet initiated by \(i\) carry fraction \(x\) of the
parent energy.  Perturbative coefficients are then convolved with \(T_i\)
rather than evaluated with a purely partonic charged-energy measurement.
The renormalization group has DGLAP form with nonlinear mixing reflecting the
fact that charged fractions of daughter branches add:
\[
  \mu\frac{\dd}{\dd\mu}T_i(x,\mu)
  =
  \sum_{j,k}\int_0^1\dd z
  \int_0^1\dd x_1\,\dd x_2\,
  \frac{\alpha_s(\mu)}{2\pi}P_{i\to jk}(z)\,
  T_j(x_1,\mu)T_k(x_2,\mu)
  \delta(x-zx_1-(1-z)x_2)
  +\cdots .
\]
The ellipsis denotes virtual terms and higher-order kernels fixed by the same
renormalization prescription.  Track functions are therefore
nonperturbative inputs with perturbative evolution, not substitutes for IRC
safety.

\section{Parton Showers}

A parton shower is a stochastic implementation of soft and collinear
factorized evolution.  At leading soft-collinear order, a final-state
branching probability for a parton of type \(i\) is built from the Sudakov
factor
\begin{equation}
  \Delta_i(t_1,t_0)
  =
  \exp\left[
    -\int_{t_0}^{t_1}\frac{\dd t}{t}
    \sum_{i\to jk}
    \int \dd z\,
    \frac{\alpha_s(\kappa(t,z))}{2\pi}
    P_{i\to jk}(z)
  \right],
  \label{eq:parton-shower-sudakov}
\end{equation}
where \(t\) is the shower ordering variable, \(z\) is the splitting fraction,
\(P_{i\to jk}\) is an Altarelli--Parisi splitting kernel, and
\(\kappa(t,z)\) is the scale at which the running coupling is evaluated.
The probability for the next branching between \(t\) and \(t+\dd t\) is the
Sudakov no-branching probability multiplied by the differential splitting
kernel.

Different showers choose different evolution variables, recoil maps, momentum
conservation prescriptions, and color structures.  Angular-ordered showers
encode color coherence through emission angles.  Dipole and antenna showers
use color-connected pairs as emitters.  Many practical implementations use a
leading-color approximation in which the color algebra is replaced by a
planar color-flow structure; subleading-color corrections require additional
operator information.

\begin{controlledapproximation}[Shower accuracy]
A shower calculation has a stated logarithmic accuracy only for a specified
class of observables and a specified matching or merging scheme.  The basic
Markov shower reproduces leading soft-collinear logarithms for observables
compatible with its ordering variable.  NLO-matched prescriptions such as
MC@NLO or POWHEG replace the first hard emission and inclusive normalization
by exact NLO matrix-element information while retaining the shower resummation
for subsequent unresolved radiation.  Multijet merging combines samples with
different resolved parton multiplicities using a merging scale whose residual
dependence is part of the perturbative uncertainty estimate.
\end{controlledapproximation}

\section{Hadronization and Fragmentation Functions}

Hadronization is the nonperturbative formation of outgoing hadrons from the
high-energy QCD state produced by the short-distance interaction.  In the
exact theory it is simply time evolution in the physical Hilbert space of
QCD.  In phenomenological calculations it is represented by additional
nonperturbative inputs after the perturbative hard, jet, and soft evolution
has been organized.

For identified hadron production, the corresponding input is a fragmentation
function.  Let \(n^\mu\) be a lightlike vector and let
\(\bar n^\mu\) satisfy \(n\cdot\bar n=2\).  A quark fragmentation function may
be represented by a light-ray operator with Wilson lines:
\begin{equation}
  D_{h/q}(z,\mu)
  =
  \frac{1}{2zN_c}
  \sum_X
  \int\frac{\dd y^+}{2\pi}
  \ee^{-\ii P_h^- y^+/z}
  \operatorname{tr}
  \left[
    \frac{\gamma\cdot\bar n}{2}
    \langle0|W^\dagger(y^+)q(y^+)|hX\rangle
    \langle hX|\bar q(0)W(0)|0\rangle
  \right]_{\rm ren}.
  \label{eq:quark-fragmentation-function-operator}
\end{equation}
Here \(P_h^-\) is the large light-cone momentum of the observed hadron,
\(z=P_h^-/k^-\) is the momentum fraction relative to the parent partonic
coordinate \(k^-\), \(W\) is the appropriate lightlike Wilson line, and the
subscript indicates operator renormalization.  The scale dependence of
\(D_{h/q}\) is governed by DGLAP evolution, while its boundary value at a
finite scale is nonperturbative data.

Lund string and cluster hadronization models are phenomenological
parametrizations of this nonperturbative stage.  A string model uses color
flux-tube breaking with fitted parameters for string tension, flavor
production, and transverse momentum.  A cluster model first organizes the
partonic state into color-singlet clusters and then decays those clusters into
hadrons with fitted rules.  Their predictions therefore carry model and tuning
uncertainties in addition to perturbative uncertainties.

The lattice formulation of
Chapter~\ref{ch:lattice-yang-mills-nonperturbative-definition} is the
regulator-level nonperturbative definition of the gauge theory used in this
volume.  It computes Euclidean finite-volume correlation functions directly in
terms of hadronic states.  Real-time fragmentation functions and fully
exclusive hadronization histories involve lightlike or timelike information,
so their extraction from Euclidean lattice data requires additional
reconstruction, finite-volume, or large-momentum matching ideas.  This is a
technical obstruction in the observable, not a license to replace the
nonperturbative QCD input by a perturbative calculation.

\section{Classes of QCD Final-State Observables}

The useful separation is by the mathematical input required by the observable.

\begin{description}[style=nextline,leftmargin=1.5em]
\item[Total inclusive rates and IRC-safe event shapes.]
The required input is renormalized fixed-order amplitudes, KLN cancellation,
and resummation when hierarchical scales are measured.

\item[IRC-safe jet cross sections.]
The required input is an IRC-safe jet algorithm, fixed-order hard functions,
PDFs for incoming hadrons, and logarithmic resummation where jet vetoes or
small radii introduce large logarithms.

\item[Energy correlators.]
The required input is the stress-tensor energy-flow operators on the physical
Hilbert space, with partonic perturbation theory used only after the detector
observable is defined.

\item[Identified hadron spectra.]
The required input is fragmentation functions or fracture functions, their
operator renormalization, and nonperturbative boundary data.

\item[Fully exclusive hadronization patterns.]
The required input is the full QCD final-state distribution or a
phenomenological model with stated fit parameters and model uncertainty.
\end{description}

This classification keeps the physics claims separated.  Perturbative QCD
computes short-distance coefficients and IRC-safe partonic representatives.
Fragmentation functions, PDFs, shape functions, and hadronization models are
additional nonperturbative data or controlled approximations with their own
definitions and uncertainties.  Energy correlators are especially valuable
because they are nonperturbatively defined by the stress tensor and are
insensitive, after smearing, to soft emissions and collinear recombination in
the precise sense used in
Proposition~\ref{prop:irc-safety-fixed-order-finiteness}.
